\section{Two filtrations and a second stabilizer}\label{filt:section}

Fix one of the gradings and points furnished by
Theorem~\ref{adj:saturation}.  We retain the original degree $j$ and put
$n=N+1$, with $N\geq4$.  Thus $S=\CC[C]$ is a positively graded normal
Gorenstein domain, its vertex $o$ is singular, and its canonical module has
a homogeneous generator of original degree $r$.  The Taylor semigroup and
its lattice satisfy
\begin{equation}\label{filt:semigroup-input}
 \Gamma=K\cap\Lambda,\qquad
 \Int(K)\cap\Lambda=c_0+\Gamma,\qquad
 c_0=(r,1,\ldots,1).
\end{equation}
Here $K$ is rational polyhedral.  We also use the finite Taylor-order
bound and the eventual realization of compatible jets from
Proposition~\ref{jet:setup}.  The only volume hypothesis in this section is
\begin{equation}\label{filt:volume-hypothesis}
 \nvol(o,C)\geq 2N^n.
\end{equation}
The first filtration separates total Taylor order from original degree.
We then assign positive cost to the constant Taylor direction as well.
The resulting second graded ring has a nontrivial stabilizer, which will
supply the same-dimensional slice in Section~\ref{boot:section}.

\subsection{The first associated graded ring}

For an original-homogeneous section $s\in S_j$, let $T(s)$ be its total
Taylor order at the chosen orbifold point, computed in the fixed smooth
uniformizing chart and local frame.  Extend the filtration by direct sums
in original degree, and write
\[
 \overline S=\gr_T S.
\]
Its bidegrees will always mean original degree and $T$-degree, in that
order.  If $\tau$ is a bookkeeping variable for original degree, then
$\overline S$ is the subalgebra of $\CC[\tau,z_1,\ldots,z_N]$ consisting of
homogeneous Taylor forms $\tau^jP(z)$ obtained from sections.  Each such
form satisfies the stabilizer congruence.  Taking its further monomial
initial with the fixed lexicographic refinement gives the algebra
$\CC[\Gamma]$.

The saturated monomial algebra will determine normality and the canonical
degree of $\overline S$.  A second consequence of saturation is a section
whose Taylor form is the constant $\tau^m$.  Inverting that section
recovers every compatible Taylor monomial; this will determine the
Hilbert series of its zero fiber and give a basis for the next filtration.

\begin{lemma}[The Taylor graded ring]\label{filt:first-ring}
The algebra $\overline S$ is finitely generated, normal, Cohen--Macaulay,
and Gorenstein.  Its canonical module has a homogeneous generator of
bidegree $(r,N)$.  There is a section $v\in S_m$ whose initial
$\overline v$ is $\tau^m$.  In particular $T(v)=0$, and
\begin{equation}\label{filt:laurent-product}
 \overline S[\overline v^{-1}]
 \simeq \CC[\overline v,\overline v^{-1},u_1,\ldots,u_N],
 \qquad T(u_i)=1,\quad T(\overline v)=0.
\end{equation}
The ring $D=\overline S/(\overline v)$ is an equidimensional
Cohen--Macaulay Gorenstein ring of dimension $N$, with
\begin{equation}\label{filt:D-series}
 \Hilb_T D(z)=(1-z)^{-N},
 \qquad \deg\omega_D=(r-m,N).
\end{equation}
Its $T$-degree-zero piece is $\CC$.  A bihomogeneous vector-space basis of
$D$ can be lifted to sections $e_i\in S$ such that, on putting $T_i=T(e_i)$,
\begin{equation}\label{filt:S-basis}
 \{v^a e_i:a\in\NN\}
\end{equation}
is a vector-space basis of $S$, compatible with original degree and
$T$-order.  The unique vector of this basis with $a+T_i=0$ is $1$.
\end{lemma}

\begin{proof}
We first transfer the structure of the monomial algebra to
$\overline S$.  By Equation~\eqref{filt:semigroup-input}, $\Gamma$ is a
finitely generated normal affine semigroup.  Hochster's theorem makes
$\CC[\Gamma]$
Cohen--Macaulay, and the canonical-module formula for normal affine
semigroup rings identifies its canonical module with the monomial ideal
of interior lattice points; see \cite{Hochster72} and
\cite[Theorem~6.7 and its following normal-monoid reformulation]{Stanley78}.
The second identity in Equation~\eqref{filt:semigroup-input} makes this
ideal principal, generated in multidegree $(r,1,\ldots,1)$.  Thus its
bidegree is $(r,N)$.  The interior-translation identity, rather than
saturation alone, is essential here.

Choose a finite set of generators of $\Gamma$ and bihomogeneous lifts in
$\overline S$ of their monomials.  We lift generation by degreewise
subduction, in the terminology of Khovanskii bases
\cite[Algorithm~2.11 and Proposition~2.13]{KavehManon2019}.
Subtracting products of the lifts cancels any selected initial monomial.
Within a fixed original degree there are only finitely many Taylor
orders, and within each order only
finitely many monomials.  Consequently this subtraction terminates and
the chosen lifts generate $\overline S$.

The same finite list gives a one-parameter flat monomial degeneration.
Choose positive integral costs on the $z_i$ that preserve the selected
monomial initial of every one of these
finitely many lifts.  Such costs exist by approximating the specified
lexicographic refinement on this finite list.  The scalar initial
algebra contains $\CC[\Gamma]$.  In each original bidegree its dimension
is the dimension of $\overline S$ in that bidegree, which is also the
number of the corresponding points of $\Gamma$.  The inclusion is
therefore an equality in every bidegree.  Lifting the generators again,
with the same degreewise terminating subtraction, shows that the
extended Rees algebra for this scalar filtration is finite type.  As a
subalgebra of a Laurent-polynomial extension it is torsion-free, hence
flat, over its parameter line.  Its zero fiber is $\CC[\Gamma]$ and its
nonzero fibers are $\overline S$.

Write $R$ for this extended Rees algebra and $\lambda$ for its parameter.
It retains original degree and $T$-degree, and has the additional scalar
filtration grading, with scalar weight $-1$ on $\lambda$.
A sufficiently large multiple of original degree minus scalar weight
gives a positive grading on $R$, including positive degree for $\lambda$.
At its vertex, the local ring $A$ has a normal Gorenstein quotient
$A/(\lambda)$.
Lifting through the nonzerodivisor $\lambda$ makes $A$
Cohen--Macaulay and Gorenstein; for the latter equivalence see
\cite[Tag~0BJJ]{StacksProject}.

The local ring $A$ is also normal.  Here is a direct verification using
$R_1+S_2$.  It is already Cohen--Macaulay.  Let $P$ be a height-one
prime of $A$.  If $\lambda\in P$, reducedness of $A/(\lambda)$ makes
$A_P/(\lambda)$ a field, so $A_P$ is regular.  Otherwise choose $Q$
minimal over $(P,\lambda)$.  Since $A$ is a local domain essentially of
finite type over $\CC$, the dimension formula and the principal ideal
theorem give $\operatorname{ht}Q\leq2$.  The normal ring
$A_Q/(\lambda)$ has dimension at most one and is regular.  Lifting a
regular system of parameters together with $\lambda$ shows that $A_Q$
is regular, and then so is its localization $A_P$.

The normal and Gorenstein loci of $\Spec R$ are open, contain its
vertex, and are invariant under the positive grading action.  Every
orbit contracts to that vertex; hence $R$ is normal and Gorenstein
everywhere.  Its generic fiber, a localization of $R$, is
$\overline S\otimes_{\CC}\CC(\lambda)$ and is normal.  Normality
descends under this faithfully flat field extension, proving normality
of $\overline S$.  Quotienting $R$ by $\lambda-1$ proves the
Cohen--Macaulay and Gorenstein assertions for $\overline S$ as well.

Adjunction tracks the canonical bidegree.  The parameter $\lambda$ has
original degree and $T$-degree both zero, so it contributes no shift in
these two degrees.  Choose a lift of the central canonical generator
homogeneous also for the scalar grading, with its adjunction shift.
Its cokernel in the canonical module of $R$ is equal to its product by
$\lambda$.  Graded Nakayama for the positive combined grading therefore
makes the cokernel zero.  Restriction to $\lambda=1$ forgets the scalar
grading but preserves original degree and $T$-degree.  Thus the canonical
generator of $\overline S$ has bidegree $(r,N)$.

We next use the constant Taylor direction to describe $\overline S$ over
a polynomial parameter.  The vector $(m,0)$ lies on the constant ray of
$K$ and belongs to $\Lambda$.  Equation~\eqref{filt:semigroup-input} provides a section with
this leading index.  Rescale its value to obtain
$\overline v=\tau^m$.  Choose integer representatives $b_i$ of the
stabilizer characters.  Eventual compatible jet separation shows that,
after inverting $\tau^m$, every monomial $\tau^jz^\gamma$ with
\[
 j-\sum_i b_i\gamma_i\in m\ZZ
\]
occurs.  Set $u_i=\tau^{b_i}z_i$ in that localization.  The displayed
congruence expresses each allowed monomial uniquely as a Laurent power
of $\overline v$ times a monomial in the $u_i$.  This proves
Equation~\eqref{filt:laurent-product}, including its $T$-grading.  It also
shows directly that there is no residual finite quotient in a nonzero
$\overline v$ fiber.

For each $q\geq0$, consider the $T$-degree-$q$ piece of $\overline S$ as
a module over $\CC[\overline v]$.  It is torsion-free.  It is also finite:
in each original-degree residue class modulo $m$, multiplication by
$\overline v$ identifies its pieces with an increasing sequence of
subspaces of the finite-dimensional space of compatible degree-$q$
Taylor forms.  These subspaces are eventually the full compatible space
by jet separation.  Thus there are only finitely many necessary module
generators.  This module is consequently free over the principal ideal
domain $\CC[\overline v]$.  Equation~\eqref{filt:laurent-product} shows
that its rank is
\[
 \binom{q+N-1}{N-1}.
\]
Reduction modulo $\overline v$ consequently has exactly this dimension
in $T$-degree $q$, proving the Hilbert series in
Equation~\eqref{filt:D-series}.

Since $\overline v$ is a nonzerodivisor in the normal
Cohen--Macaulay domain $\overline S$, its quotient is equidimensional of
dimension $N$ and Cohen--Macaulay.  Adjunction makes the quotient
Gorenstein and subtracts $(m,0)$ from the canonical degree.  The Hilbert
series shows that the $T$-degree-zero piece of $D$ is $\CC$.

It remains to obtain a basis of $S$ compatible with the filtration.
Lift a bihomogeneous basis of $D$ to $\overline S$.  These lifts are a
$\CC[\overline v]$-basis: generation follows by subtracting modulo
$\overline v$ and descending in original degree, and independence follows
by reducing a putative relation modulo $\overline v$ and successively
removing its smallest power.  Now lift them through the $T$-filtration
to original-homogeneous sections $e_i$ of the indicated orders $T_i$.
Their products with powers of $v$ have the just-constructed basis as
initials.  Terminating cancellation in original degree proves
Equation~\eqref{filt:S-basis}.  Choose $e_i=1$ for the sole basis vector
of $T$-degree zero; all other $T_i$ are positive.
\end{proof}

\subsection{Canonical forms on extended Rees spaces}

We make the discrepancy comparison used for both filtrations explicit.
If $F$ is a nonnegative decreasing multiplicative integral filtration on
$S$, set
$F_pS=S$ for $p\leq0$ and use the extended Rees convention
\begin{equation}\label{filt:rees-convention}
 R_F=\sum_{p\in\ZZ}\lambda^{-p}F_pS
       \subset S[\lambda,\lambda^{-1}].
\end{equation}
The filtration action has weight zero on $S$ and weight $-1$ on
$\lambda$.  In particular, $R_F/(\lambda)=\gr_F S$ and
$R_F[\lambda^{-1}]=S[\lambda,\lambda^{-1}]$.  The inclusion
$S[\lambda]\subset R_F$ gives a birational morphism
$\Spec R_F\to C\times\A^1$.

\begin{lemma}[Rees discrepancy comparison]\label{filt:discrepancy}
Suppose $F$ respects original degree, has finite length in each original
homogeneous piece, and $\gr_F S$ is finitely generated.  Suppose also
that $\mathfrak X=\Spec R_F$ is normal and Gorenstein, and that
$\gr_F S$ is Cohen--Macaulay and Gorenstein with a canonical generator
of bidegree $(r,Q)$.  Then $R_F$ is finite type and its absolute canonical
module has a homogeneous generator of bidegree $(r,Q-1)$, which on
$\lambda\ne0$ is
\begin{equation}\label{filt:canonical-form}
 \Omega_{\mathfrak X}
   =c\lambda^{-Q}\Omega_C\wedge d\lambda,
 \qquad c\in\CC^*.
\end{equation}
For a quasi-monomial valuation $V$ centered on $\mathfrak X$ over
$\lambda=0$, put $\ell=V(\lambda)>0$ and
$w=V|_{\CC(C)}$.  If $w$ is centered at $o$, then
\begin{equation}\label{filt:discrepancy-inequality}
 A_C(w)\leq A_{\mathfrak X}(V)+(Q-1)\ell.
\end{equation}
The same inequality holds whenever the discrepancies involved are
finite and defined by valuation retractions.
\end{lemma}

\begin{proof}
Lift homogeneous generators of $\gr_F S$ to $s_i\in S$, of filtration
orders $p_i$.  Cancellation of initial forms, terminating within each
original degree, shows that $R_F$ is generated by $\lambda$ and
$\lambda^{-p_i}s_i$.  Indeed, an element of $F_pS$ has an expression as
a sum of monomials in the $s_i$ whose filtration orders are at least $p$;
multiplication by $\lambda^{-p}$ then introduces only nonnegative extra
powers of $\lambda$.  Taking a sufficiently large multiple of original
degree minus filtration degree makes every one of these generators
positive, with $\deg\lambda=1$.

Cartier adjunction \cite[Tag~0B4A]{StacksProject} for the nonzerodivisor
$\lambda$ identifies the reduction of
the canonical module of $R_F$ with that of $\gr_FS$, with the shift
coming from $d\lambda$.  Since $\lambda$ has filtration weight $-1$, a
lift of the specified central generator has bidegree $(r,Q-1)$.
This lift generates the whole canonical module: its cokernel is equal
to its product by $\lambda$, and a finite module bounded below in the
positive combined grading cannot have that property unless it is zero.
As the canonical module is torsion-free of rank one over the domain
$R_F$, the resulting generator is free.

On $\lambda\ne0$, both this generator and
$\Omega_C\wedge d\lambda$ generate the same canonical module.  Their
ratio is a unit of $S[\lambda,\lambda^{-1}]$.  The units of the
positively graded domain $S$ are the nonzero constants, so such a unit
is $c\lambda^a$.  The filtration weight of
$\Omega_C\wedge d\lambda$ is $-1$, and the weight of the lifted generator
is $Q-1$.  It follows that $a=-Q$, proving
Equation~\eqref{filt:canonical-form}.  Computing on a common smooth
birational model gives
\begin{equation}\label{filt:discrepancy-product}
 A_{C\times\A^1}(V)
   =A_{\mathfrak X}(V)+Q\ell.
\end{equation}
This computation uses rational top forms and does not require the
Rees morphism itself to be proper.

On the other hand,
\begin{equation}\label{filt:product-lower}
 A_{C\times\A^1}(V)\geq A_C(w)+\ell.
\end{equation}
To see this directly, take any smooth proper birational log model of
$C$ used to retract $w$, and take its product with the parameter line.
Include $\lambda=0$ among the simple-normal-crossing divisors.  The
monomial retraction of $V$ on this product has discrepancy equal to the
discrepancy of the corresponding retraction of $w$ plus $\ell$.
Log discrepancy dominates its value on each such retraction
\cite[Corollary~5.4 and Remark~5.6]{JM12}; taking
the supremum over the models of $C$ proves
Equation~\eqref{filt:product-lower}.  Combining it with
Equation~\eqref{filt:discrepancy-product} proves the lemma.
\end{proof}

\subsection{Reducedness and the second filtration}

\begin{lemma}[Reducedness of the quotient]\label{filt:reduced}
Under Equation~\eqref{filt:volume-hypothesis}, $D$ is reduced.
\end{lemma}

\begin{proof}
The $T$-Rees algebra is finite type by Lemma~\ref{filt:first-ring} and
the generator-lifting argument in Lemma~\ref{filt:discrepancy}.
Its nonzero fibers are $S$, and its central fiber $\overline S$ is
normal, Cohen--Macaulay, and Gorenstein.  The regular parameter and these
fiber properties imply that the total space is Cohen--Macaulay and
Gorenstein.  It is normal as well: away from the central fiber it is
$C\times\Gm$, and at every codimension-one point of the central fiber
the quotient by $\lambda$ is a field.  Thus it is regular in codimension
one and satisfies Serre's condition $S_2$.

Let $H$ be an irreducible component of the divisor $\overline v=0$ in
$\Spec\overline S$, and write
$\nu=\ord_H(\overline v)\geq1$.  At its generic point the normal central
fiber is a discrete valuation ring.  The total space is therefore
regular there, of local dimension two, with parameters $\lambda$ and a
lift $t$ of a central uniformizer.  In this local ring
\[
 v=a t^{\nu}+\lambda b
\]
for a unit $a$ and a regular element $b$.  Take the monomial valuation
with costs $V(\lambda)=1$ and $V(t)=1/\nu$.  Its discrepancy on the
total space is $1+1/\nu$.  For its restriction $w$ to $\CC(C)$ one has
\[
 w(v)\geq1,\qquad w(e_i)\geq T_i,
\]
because $\lambda^{-T_i}e_i$ is regular on the Rees space.  Hence
$w(v^ae_i)\geq a+T_i$.

By Lemma~\ref{filt:first-ring}, every nonconstant original-homogeneous
basis vector has a strictly positive lower bound $a+T_i$.  Thus $w$ is
centered exactly at $o$.  The series counting the basis with these
costs is
\[
 \sum_{a,i}z^{a+T_i}=(1-z)^{-1}\Hilb_TD(z)=(1-z)^{-n}.
\]
The vectors of cost below $k$ consequently number
$k^n/n!+O(k^{n-1})$.  They span the valuation quotient, so
\begin{equation}\label{filt:first-volume}
 \vol(w)\leq1.
\end{equation}
These finite-colength quotients of $S$ are supported at $o$; their
lengths agree with those after localization at the vertex.

Apply Lemma~\ref{filt:discrepancy} with $Q=N$ and $\ell=1$.  It gives
\[
 A_C(w)\leq N+1/\nu.
\]
If $\nu\geq2$, Equation~\eqref{filt:first-volume} yields
\[
 \nvol(o,C)\leq (N+1/2)^{N+1}<2N^{N+1}.
\]
For the strict inequality, use
\[
 (N+1)\log\left(1+\frac1{2N}\right)
 <\frac{N+1}{2N}\leq\frac58<\log2,
 \qquad N\geq4.
\]
This contradicts Equation~\eqref{filt:volume-hypothesis}.  All the
multiplicities of $\overline v$ are therefore one.  Its quotient $D$ is
generically reduced, and it is Cohen--Macaulay, hence satisfies $S_1$.
The reducedness criterion $R_0+S_1$ proves the assertion; see
\cite[Tag~033P]{StacksProject}.
\end{proof}

Define a second filtration by assigning the basis in
Equation~\eqref{filt:S-basis} the costs
\begin{equation}\label{filt:W-definition}
 W(v^a e_i)=a+T_i
\end{equation}
and taking spans of basis vectors of cost at least the indicated level.
Unlike $T$, this filtration is positive on every nonconstant
original-homogeneous vector.

There is also a description independent of the lifted basis. Write
$T_qS=\{s:T(s)\ge q\}$, with $T_qS=S$ for $q\le0$.
For the distinguished section $v$ fixed above, one has
\begin{equation}\label{filt:W-intrinsic}
 W_pS=\sum_{a=0}^{p}v^aT_{p-a}S\qquad(p\ge0).
\end{equation}
Indeed, every basis term $v^be_i$ on the right has
$b+T_i\ge p$. Conversely, for such a term choose
$a=\min\{p,b\}$; then $v^{b-a}e_i\in T_{p-a}S$.
Thus changing the compatible lifts $e_i$ does not change $W$.
The multiplicativity of $T$ also makes the right side of
Equation~\eqref{filt:W-intrinsic} multiplicative. The remaining
point is to identify the associated graded algebra, not merely
its Hilbert function.

\begin{proposition}[The two-filtration structure]\label{filt:structure}
The filtration $W$ is multiplicative, with
\begin{equation}\label{filt:second-graded}
 G=\gr_WS\simeq D[\mathtt v],\qquad
 \deg_W\mathtt v=1,\qquad \deg_W|_D=T.
\end{equation}
Here $\mathtt v$ is a regular polynomial variable of original degree
$m$.  Moreover
\begin{equation}\label{filt:G-series}
 \Hilb_WG(z)=(1-z)^{-n},\qquad
 \deg\omega_G=(r,n).
\end{equation}
The $W$-Rees space is a finite-type normal Gorenstein variety.  Its
canonical generator has bidegree $(r,n-1)$ and, on $\lambda\ne0$, is
$c\lambda^{-n}\Omega_C\wedge d\lambda$.
\end{proposition}

\begin{proof}
Expand a product $e_ie_j$ in the basis $v^ae_k$.  Compatibility of that
basis with total Taylor initials shows that every nonzero term has
$T_k\geq T_i+T_j$.  Its $W$-cost $a+T_k$ is therefore at least
$T_i+T_j$.  Terms with equality have precisely $a=0$ and
$T_k=T_i+T_j$; their images are the product of the images of $e_i,e_j$
in $D=\overline S/(\overline v)$.  Factoring out the displayed powers
of $v$ proves multiplicativity and the algebra isomorphism in
Equation~\eqref{filt:second-graded}.  In particular the variable
$\mathtt v$ is regular even if $D$ has several components.

The Hilbert series follows from Equation~\eqref{filt:D-series}.
Adjoining the polynomial variable of bidegree $(m,1)$ changes the
canonical degree $(r-m,N)$ to $(r,n)$.  Thus $G$ is reduced,
Cohen--Macaulay, and Gorenstein.  The $W$-filtration has finite length in
fixed original degree, as is clear from its basis description; lifting
finitely many generators of $G$ proves finite generation of its Rees
algebra.

The nonzero fibers of this Rees family are $S$, and its central fiber
is $G$.  Consequently the total space is Cohen--Macaulay and Gorenstein.
For normality, away from $\lambda=0$ it is $C\times\Gm$.  A
codimension-one point on $\lambda=0$ is the generic point of a central
component.  Since $G$ is reduced, its local ring there is a field, and
the total local ring has maximal ideal generated by the nonzerodivisor
$\lambda$.  It is regular.  The total space satisfies $R_1$ and $S_2$,
hence is normal.  Notice that normality of $D$ or of $G$ is not required
in this argument.  Finally Lemma~\ref{filt:discrepancy}, with $Q=n$,
gives the canonical generator and top-form identity.
\end{proof}

\Needspace{12\baselineskip}
\paragraph{Example: the quadratic cone.}
The two stabilizers and filtrations can already differ for an ordinary
double point.  Let
\[
 S=\CC[a,z_1,\ldots,z_N,b]/(ab-q(z)),
 \qquad q(z)=\sum_{i=1}^N z_i^2,
\]
with all original degrees equal to one.  Here \(r=N\) and \(m=1\).
At the quotient point \([a:z:b]=[1:0:0]\), the chart \(a=1\)
has \(b=q(z)\).  Taking \(v=a\), the three degrees are
\[
\begin{array}{c@{\qquad}ccc}
 \toprule
 &a&z_i&b\\
 \midrule
 \text{original degree}&1&1&1\\
 T&0&1&2\\
 W&1&1&2\\
 \bottomrule
\end{array}
\]
The first Taylor graded ring has the same presentation as \(S\), and
\(D=\CC[z_1,\ldots,z_N,b]/(q(z))\).  The family
\(\Spec\overline S\to\A^1_t\) defined by \(t=\overline v=a\)
has invariant parameter for the \(T\)-action.  After inverting \(t\),
\(b=q(z)/a\) gives the polynomial family of
Equation~\eqref{filt:laurent-product}.  With
\(\mathsf a=\lambda^{-1}a\), \(\mathsf z_i=\lambda^{-1}z_i\),
and \(\mathsf b=\lambda^{-2}b\), the second Rees space is
\[
 \mathfrak X=
 \{\lambda\mathsf a\mathsf b-q(\mathsf z)=0\}.
\]
Its action has weights \((-1,1,1,2)\) on
\((\lambda,\mathsf a,\mathsf z_i,\mathsf b)\).
The central point with \(\mathsf b=1\) and all other coordinates zero
therefore has stabilizer \(h=2\).  Its transverse slice
\[
 Y=\{\mathsf b=1\}
   =\{\lambda\mathsf a-q(\mathsf z)=0\}
\]
is again an ordinary double point of dimension \(n=N+1\).
Thus the original stabilizer \(m=1\) and the new stabilizer \(h=2\)
play different roles, and taking this slice preserves the original
dimension.

\subsection{A nontrivial stabilizer on the new central ring}

\begin{lemma}[Free gradings and polynomial rings]\label{filt:free-grading}
Let $E$ be a finitely generated positively graded reduced
equidimensional Cohen--Macaulay $\CC$-algebra of dimension $N\geq2$,
with $E_0=\CC$ and $\Hilb E(z)=(1-z)^{-N}$.  If the grading action has
trivial stabilizer at every nonvertex closed point of $\Spec E$, then
$E$ is a polynomial ring on $N$ generators of degree one.
\end{lemma}

\begin{proof}
We first justify that weight one defines an honest line bundle on
$Z=\Proj E$.  Choose positive homogeneous algebra generators $a_j$,
with degrees $w_j$.  At a nonvertex closed point $q$, its grading
stabilizer is
\[
 \mu_{\gcd\{w_j:a_j(q)\ne0\}}.
\]
By hypothesis this gcd is one.  Invert the product $f$ of those active
generators.  B\'ezout's identity supplies integers $c_j$ with
$\sum_j c_jw_j=1$, so $u=\prod_j a_j^{c_j}$ is an invertible homogeneous
element of degree one in $E_f$.  Consequently
\begin{equation}\label{filt:torsor-chart}
 E_f=(E_f)_0[u,u^{-1}].
\end{equation}
These open sets cover the puncture: they contain all its closed points,
and a finite-type scheme over $\CC$ is Jacobson.  Their degree-zero
charts cover $Z$.  Equation~\eqref{filt:torsor-chart} gives a Zariski
$\Gm$-torsor over $Z$, and makes $L_Z=\cO_Z(1)$ invertible, with
$\cO_Z(k)=L_Z^{\otimes k}$.  A positive power of this line bundle is
ample by the construction of $\Proj$, so $L_Z$ is ample.

The same charts show that $Z$ is reduced, equidimensional, and
Cohen--Macaulay, of dimension $N-1$.  The Hilbert leading coefficient
gives $L_Z^{N-1}=1$.  Each irreducible projective component has a
positive integral degree for this actual ample line bundle.  There is
therefore only one component, and reducedness makes $Z$ integral.
The image $V_0$ of $E_1$ in $H^0(Z,L_Z)$ has dimension $N$: restriction
is injective, since a nonzero homogeneous element in a reduced
positive-dimensional equidimensional ring cannot vanish on its whole
puncture.

We construct successive Cartier sections while tracking the images of
$V_0$, not assuming surjectivity onto full spaces of sections.  Suppose
$Z_i$ is a positive-dimensional integral projective Cohen--Macaulay
scheme obtained after $i$ cuts, and that the image $V_i$ of $V_0$ on it
has dimension $N-i$.  Choose a nonzero section $s_i\in V_i$.  Its zero
scheme $Z_{i+1}$ is an effective Cartier divisor.  It is nonempty,
because $L_Z|_{Z_i}$ is ample of positive degree, and it is pure and
Cohen--Macaulay.  Its degree is again one.  The component-degree formula
forces a single component of generic multiplicity one; the $S_1$
property then makes this divisor reduced, hence integral.  The kernel
of restriction $V_i\to H^0(Z_{i+1},L_Z|_{Z_{i+1}})$ is exactly
$\CC s_i$: division by $s_i$ identifies that kernel with regular global
functions on the proper integral scheme $Z_i$, and those functions are
constant.  Thus $\dim V_{i+1}=N-i-1$.

After $N-1$ cuts, the resulting scheme is a single reduced point and
the remaining image space has dimension one.  Choose a final section
nonzero at that point.  Lifting all these sections to $E_1$ gives $N$
elements whose common zero on $\Proj E$ is empty.  Their ideal is
therefore irrelevant-primary and they form a homogeneous system of
parameters.  Cohen--Macaulayness makes the sequence regular.  The
quotient has Hilbert series
\[
 (1-z)^N\Hilb E(z)=1.
\]
It is $\CC$, so induction on degree, or graded Nakayama, shows that
these $N$ linear elements generate $E$.  The resulting surjection from
a polynomial ring on $N$ degree-one variables is an isomorphism by the
same Hilbert series.
\end{proof}

\begin{corollary}[A nontrivial second stabilizer]\label{filt:stabilizer}
There is a nonvertex closed point of $\Spec D$ whose $T$-grading
stabilizer has order $h\geq2$.
\end{corollary}

\begin{proof}
Otherwise Lemma~\ref{filt:free-grading} would make $D$ a polynomial ring
on $N$ degree-one generators, and Proposition~\ref{filt:structure} would
make $G$ a polynomial ring on $n=N+1$ generators.  Choose these generators
homogeneous also for original degree and lift them to $S$.  Terminating
cancellation in the $W$-filtration shows that the lifts generate $S$.
A surjection from a polynomial ring on $n$ variables to the
$n$-dimensional domain $S$ has zero kernel.  Thus $C$ would be smooth,
contrary to our standing hypothesis.  Stabilizers at nonvertex points
of a positive grading are finite, so the asserted order $h\geq2$
exists.
\end{proof}
